\documentclass[12pt]{article}
\usepackage[T1]{fontenc}
\usepackage[utf8]{inputenc}
\usepackage[brazil]{babel}
\usepackage[margin=1in]{geometry}
\setlength{\parindent}{0pt}
\begin{document}
\section*{Tema 69}
Processo: PEDILEF 2007.71.52.004219-0/ RS\\
Situacao: Julgado\\
Ramo: DIREITO ADMINISTRATIVO\\
Questao: Saber qual o fator de divisão para o cálculo da hora extra para o servidor público.\\
Tese: O fator de divisão para o cálculo do serviço extraordinário é de 200 horas mensais, em razão da jornada máxima do servidor público de 40 (quarenta) horas semanais prevista no art. 19 da Lei n. 8.112/90.\\
Fonte: https://www.cjf.jus.br/phpdoc/virtus/
\end{document}
